\chapter{جمع‌بندی و راهکار‌های آتی}

در این پژوهش هدف تمرکز بر نقطه‌ی ضعف اصلی مدل‌های پخشی یعنی سرعت کم نمونه‌برداری آنها بود. در وهله‌ی اول با مطالعه و دسته‌بندی روش‌های موجود برای افزایش سرعت به یک دسته‌بندی از روش‌های کاهش تعداد گام در مقابل روش‌های کاهش محاسبات هر گام رسیدیم و سعی کردیم در هر دسته تعدادی روش را که بدون نیاز به آموزش بودند مورد مطالعه قرار دهیم. این مطالعات ما را به روشی تحت عنوان ادغام توکن‌ها رساند.
\\
روش ادغام توکن‌ها با در نظر گرفتن اینکه هزینه‌ی محاسباتی بلوک‌های مبدل داخل شبکه‌ی یو زیاد است و اینکه خود مبدل هزینه‌ی محاسباتی درجه دو نسبت به اندازه‌ی ورودی دارد، تلاش می‌کند تا با ترکیب توکن‌های ورودی مبدل از تعداد آنها بکاهد. روش اولیه‌ی ارائه شده در این رویکرد توکن‌های ورودی را به دو دسته‌ی مبدا و مقصد تقسیم می‌کند و سپس برای هر توکن مبدا مشابه‌ترین توکن مقصد را در نظر می‌گیرد و به این ترتیب یک گراف دو بخشی می‌سازد. سپس $r$ درصد از مشابه‌ترین توکن‌های مبدا و مقصد را با عملگر میانگین‌گیری ترکیب می‌کند.
\\
این روش نقاط ضعفی از جمله در نظر گرفتن یک درصد ثابت برای ترکیب توکن‌ها و نیز محاسبه‌ی ماتریس مشابهت در آغاز هر مبدل داشت که خود دارای هزینه‌ی محاسباتی درجه دو است. از همین رو تلاش این پژوهش بر این بوده است که این نقاط ضعف را بررسی و برای آنها راه‌حلی ارائه دهد. مشکل اول که ثابت بودن درصد ادغام در همه‌ی بلوک‌ها و همه‌ی بازه‌های زمانی و در نتیجه‌ی آن انعطاف ناپذیری روش ادغام توکن‌هاست را با ارائه‌ی راه‌حلی مبنی بر در نظر گرفتن یک آستانه‌ی $t$ و ادغام توکن‌هایی که مشابهت بیشتر از $t$ دارند و تطبیقی کردن رویکرد ادغام توکن‌ها سعی در برطرف ساختن آن داشتیم.
\\
و همچنین برای کاهش حجم سربار محاسباتی روش ادغام توکن‌ها از یک روش مبتنی بر حافظه‌ی پنهان برای ذخیره‌ی محاسبات در تعدادی از گام‌ها و فراخوانی مقادیر این محاسبات از حافظه‌ی پنهان در سایر گام‌ها استفاده کردیم. این روش بر مبنای این مشاهده در قسمت‌های مختلف شبکه‌ی یو است که اشاره می‌کند این تغییرات خروجی بلوک‌های مختلف در گام‌های زمانی متوالی ناچیز و هموار است. موضوعی که با رسم فاصله‌ی جاکارد جفت‌های ادغام شده در گام‌های متوالی در ورودی مبدل نیز قابل مشاهده است.
\\
نتایج اعمال این دو راه‌حل روش پیشنهادی ما تحت عنوان ادغام توکن تطبیقی مبتنی بر حافظه‌ی پنهان را می‌سازد. این روش بدون کاهش کیفیت می‌تواند کاهش زمانی بیشتری نسبت به روش ادغام توکن‌ها ارائه دهد. همچنین با توجه به اینکه حداکثر کاهش در روش ادغام توکن‌ها زمانی رخ می‌دهد که همه‌ی توکن‌های مبدا و مقصد را ادغام کنیم، می‌توان نشان داد که روش پیشنهادی ما این حداکثر کاهش زمانی را بدون از دست دادن کیفیت به دست می‌آورد.
\\


در ادامه برای بهبود و کارایی این روش پیشنهاداتی قابل ارائه است که به شرح زیر می‌باشند:
\begin{itemize}
	\item استفاده از روش پیشنهادی در کنار روش استفاده از حافظه‌ی پنهان عمیق\LTRfootnote{Deep Cache}: روش پیشنهادی ما تنها بر بلوک‌های با ویژگی سطح پایین و یا بلوک‌های با مقدار  $DownSample\ =\ 1$ اعمال می‌شود و باعث کاهش حجم محاسبات در این بلوک‌ها می‌شود. از طرفی روش استفاده از حافظه‌ی پنهان عمیق بلوک‌های میانی را در حافظه‌ی پنهان ذخیره می‌کند. ترکیب این دو روش باعث می‌شود که در تمامی بلوک‌ها یک رویکرد برای کاهش محاسبات اعمال شده باشد. البته ممکن است نیاز باشد این دو روش با روش سومی نیز برای افزایش کیفیت ترکیب شوند تا کیفیت در این فرآیند کاهش محاسبات افت نکند.
	
	\item اعمال روش پیشنهادی بر سایر معماری‌های موجود در مدل‌های پخشی: ساختار و معماری در نظر گرفته شده در سراسر این آزمایش مدل پخشی مبتنی بر شبکه‌ی یو است، حال آنکه اخیرا دسته‌های دیگری از مدل‌های پخشی از جمله مدل‌های پخشی مبتنی بر مبدل ارائه شده‌اند. بررسی عملکرد روش پیشنهادی و کلا رویکرد‌های کاهش توکن در این معماری‌ها موضوعی  مغفول مانده از سمت محققین این حوزه است که بررسی امکان و تنظیم این روش برای سایر معماری‌ها می‌تواند یکی از راهکار‌های پیش رو باشد.
	
	\item ترکیب روش پیشنهادی با سایر روش‌های عمود بر آن: در راهکار اول یک نمونه‌ی خاص از ترکیب روش پیشنهادی با روش‌های دیگر را بررسی کردیم. اما در این راهکار با توجه به قرار داشتن هر دو روش در یک دسته امکان افت کیفیت وجود داشت در حالیکه ترکیب روش پیشنهادی با سایر روش‌های کاهش گام که بدون کاهش کیفیت این عمل را انجام می‌دهند می‌تواند رویکرد دیگری باشد. استفاده از تقطیر دانش و ساخت شبکه‌ی دانش‌آموز و کاهش گام و یا استفاده از حل‌کننده‌های مختلف از جمله رویکرد‌هایی هستند که می‌توانند عمود بر رویکرد ارائه شده پیشنهاد شوند.
	
	\item ارائه‌ی روش‌های جدید ادغام یا رفع ادغام: در معرفی انواع روش‌های ادغام در فصل دوم اشاره شد که از رویکرد‌هایی مثل خوشه‌بندی نیز می‌توانیم برای ادغام توکن‌ها استفاده کنیم. علت عمده‌ی عدم استفاده از این روش‌ها زمانبر بودن آنهاست. در روش پیشنهادی با توجه به استفاده از ایده‌ی به کارگیری حافظه‌ی پنهان می‌توان در خصوص سایر روش‌های ادغام مانند خوشه‌بندی یا رویکرد‌های مشابه هم فکر کرد. و در نهایت در روش پیشنهادی رفع ادغام از طریق جایگذاری توکن مقصد در تمامی توکت‌های مبدا ادغام شده با آن می‌باشد این رویکرد نیز می‌تواند بهبود پیدا کند و از روش‌های پیچیده‌تری برای ادغام استفاده شود.
	
\end{itemize}










